\begin{mistake}{2026-04-21}{05}

\begin{problemcontent}
设 \(\alpha_1 = \begin{bmatrix} 4 \\ 1 \\ 2 \end{bmatrix}\), \(\alpha_2 = \begin{bmatrix} 1 \\ 1 \\ -1 \end{bmatrix}\), \(\alpha_3 = \begin{bmatrix} -2 \\ 0 \\ 4 \end{bmatrix}\), \(\alpha_4 = \begin{bmatrix} 7 \\ 2 \\ -3 \end{bmatrix}\), \(A = [\alpha_1, \alpha_2, \alpha_3, \alpha_4]\)，则 \(Ax = \alpha_2 + 3\alpha_4\) 的通解为\underline{\hspace{2cm}}。
\end{problemcontent}

\begin{solutioncontent}
由矩阵乘法定义，\(Ax = x_1\alpha_1 + x_2\alpha_2 + x_3\alpha_3 + x_4\alpha_4\)。
对照方程 \(Ax = \alpha_2 + 3\alpha_4\)，可直接观察出方程的一个特解为：
\[
\xi^* = \begin{bmatrix} 0 \\ 1 \\ 0 \\ 3 \end{bmatrix}
\]
接下来对系数矩阵 \(A\) 进行初等行变换，求对应齐次方程组 \(Ax=0\) 的基础解系：
\[
\begin{aligned}
A &= \begin{bmatrix} 4 & 1 & -2 & 7 \\ 1 & 1 & 0 & 2 \\ 2 & -1 & 4 & -3 \end{bmatrix} \xrightarrow{r_1 \leftrightarrow r_2} \begin{bmatrix} 1 & 1 & 0 & 2 \\ 4 & 1 & -2 & 7 \\ 2 & -1 & 4 & -3 \end{bmatrix} \\
&\to \begin{bmatrix} 1 & 1 & 0 & 2 \\ 0 & -3 & -2 & -1 \\ 0 & -3 & 4 & -7 \end{bmatrix} \xrightarrow{r_3 - r_2} \begin{bmatrix} 1 & 1 & 0 & 2 \\ 0 & -3 & -2 & -1 \\ 0 & 0 & 6 & -6 \end{bmatrix} \\
&\to \begin{bmatrix} 1 & 1 & 0 & 2 \\ 0 & -3 & -2 & -1 \\ 0 & 0 & 1 & -1 \end{bmatrix} \to \begin{bmatrix} 1 & 0 & 0 & 1 \\ 0 & 1 & 0 & 1 \\ 0 & 0 & 1 & -1 \end{bmatrix}
\end{aligned}
\]
对应的同解齐次方程组为：
\[
\begin{cases} x_1 + x_4 = 0 \\ x_2 + x_4 = 0 \\ x_3 - x_4 = 0 \end{cases}
\]
令自由变量 \(x_4 = 1\)，得基础解系 \(\eta = \begin{bmatrix} -1 \\ -1 \\ 1 \\ 1 \end{bmatrix}\)。
由非齐次方程组通解结构定理，原方程的通解为：
\[
x = k\eta + \xi^* = k \begin{bmatrix} -1 \\ -1 \\ 1 \\ 1 \end{bmatrix} + \begin{bmatrix} 0 \\ 1 \\ 0 \\ 3 \end{bmatrix} \quad (k \text{ 为任意常数})
\]
\end{solutioncontent}

\mistakeknowledge{
\begin{itemize}
 \item \textbf{核心公式/定理}：非齐次方程组通解结构 \(x = k\eta + \xi^*\)；矩阵乘法的列向量展开式 \(Ax = \sum_{i=1}^n x_i\alpha_i\)。
 \item \textbf{易错点}：遇到 \(Ax = \beta\) 右端项是特定列向量线性组合时，未直接拼凑特解而硬算增广矩阵导致计算量大易错。
 \item \textbf{关联知识}：初等行变换不改变矩阵的零空间（齐次方程组的解）；提取基础解系时，注意移项变号。
\end{itemize}}

\end{mistake}
